\chapter{Data Integration and Connectivity}
\index{on-premises data gateway}\index{VNet data gateway}\index{hybrid connectivity}\index{SQL Server!on-premises}\index{Oracle}\index{SAP}\index{REST API ingestion}\index{webhooks}\index{gateway cluster}%

\begin{objectives}
\begin{itemize}
    \item Distinguish the on-premises data gateway from the VNet data gateway and select the correct one per scenario.
    \item Connect Fabric to hybrid sources including SQL Server, Oracle, and SAP.
    \item Ingest REST APIs and webhooks with sound credential and paging patterns.
    \item Install and configure an on-premises data gateway and run a daily extraction pipeline through it.
    \item Design gateway clusters for high availability and load balancing.
    \item Architect a secure hybrid connectivity pattern for a bank that forbids public endpoints on transactional databases.
\end{itemize}
\end{objectives}

\begin{crosscloud}
Hybrid connectivity---reaching data behind corporate networks and inside VNets without exposing it publicly---is a universal enterprise requirement with remarkably similar shapes across clouds.

\vspace{2mm}
{\footnotesize
\begin{tabular}{@{}p{2.8cm}p{3.2cm}p{3.2cm}p{3.2cm}p{3.2cm}@{}}
\toprule
\textbf{Capability} & \textbf{Fabric} & \textbf{AWS} & \textbf{GCP} & \textbf{Snowflake} \\
\midrule
On-prem relay/agent & On-premises data gateway & DMS agents, DataSync agents, Storage Gateway & Transfer appliance / Database Migration Service & External access via connectors, Snowpipe through private links \\
VNet/VPC private path & VNet data gateway & VPC endpoints, PrivateLink & Private Service Connect, VPC peering & PrivateLink / Private Service Connect \\
Outbound-only pattern & Gateway initiates outbound 443 & Agents poll outbound & Agents/connectors outbound & Outbound from customer network \\
HA model & Gateway clusters (2+ nodes) & Multi-AZ services & Regional redundancy & Multi-AZ by default \\
API ingestion & Copy/Web activities, eventstreams & API Gateway + Lambda, AppFlow & Cloud Functions / Workflows & External functions, connectors \\
\bottomrule
\end{tabular}}

\vspace{1mm}
Databricks reaches private sources through VPC peering and PrivateLink; MongoDB Atlas through private endpoints and VPC peering. The invariant across clouds: \emph{the private side initiates the connection outbound}, so no inbound firewall rule ever exposes the database.
\end{crosscloud}

\begin{keyterms}
\textbf{On-premises data gateway}: an agent installed inside the corporate network that brokers outbound-only connections between Fabric and local sources. \quad
\textbf{VNet data gateway}: a Microsoft-managed gateway that lets Fabric reach sources inside an Azure virtual network without a customer-managed VM. \quad
\textbf{Gateway cluster}: two or more gateway nodes providing high availability and load distribution. \quad
\textbf{Webhook}: an HTTP callback a source system invokes to push events, rather than waiting to be polled. \quad
\textbf{Private endpoint}: a network interface that gives a service a private IP inside a VNet, removing public exposure.
\end{keyterms}

\section{Week 11 Roadmap}
Week~11 connects Fabric to the enterprise as it actually exists: databases behind firewalls, services inside VNets, and APIs that push as well as pull. Monday compares the on-premises data gateway with the VNet data gateway. Tuesday connects the classic hybrid sources---SQL Server, Oracle, SAP. Wednesday covers REST APIs and webhooks. Thursday is the hands-on project: install the on-premises gateway against a mock SQL Server and run a daily extraction to OneLake. Friday architects the bank scenario: strict hybrid connectivity with zero public endpoints.

Connectivity is where elegant cloud architectures meet organizational reality. The pipelines of Chapter~9 assume reachability; this week supplies it, securely, with the operational discipline (clustering, patching, monitoring) that makes a gateway production infrastructure rather than a developer's side install \cite{microsoftOnPremGateway,microsoftVNetGateway}.

\begin{notebox}
Gateways are identity-bearing infrastructure: the gateway's registered connections carry credentials to your most sensitive systems. Treat a gateway server with the same change control as a database server, because that is effectively what it fronts.
\end{notebox}

\section{Monday: On-Premises versus VNet Data Gateways}
\index{on-premises data gateway!versus VNet}
Both gateway types solve the same problem---Fabric is a SaaS platform; your data is not---but for different network topologies \cite{microsoftOnPremGateway,microsoftVNetGateway}.

\begin{table}[H]
\centering
\small
\begin{tabular}{@{}p{3.6cm}p{5.0cm}p{5.0cm}@{}}
\toprule
\textbf{Aspect} & \textbf{On-premises data gateway} & \textbf{VNet data gateway} \\
\midrule
Where it runs & Customer-managed Windows machine inside the corporate network & Microsoft-managed compute attached to your Azure VNet \\
Reaches & SQL Server, Oracle, file shares, SAP---anything the host can route to & Azure services inside the VNet (Azure SQL MI, VMs, private endpoints) \\
Management burden & You patch, monitor, and cluster the host(s) & Microsoft manages the nodes; you manage the VNet attachment \\
Best for & True on-premises estates, non-Azure networks & Azure-resident private sources; no VM to maintain \\
\bottomrule
\end{tabular}
\caption{Choosing between the two gateway types.}
\label{tab:chapter11-gateways}
\end{table}

\begin{definitionbox}
An \textbf{on-premises data gateway} is a customer-installed agent that creates outbound-only, encrypted connections to the cloud relay, so cloud services can query on-premises sources without any inbound firewall opening.
\end{definitionbox}

The outbound-only property is the security cornerstone: the gateway initiates every connection over TLS to the Azure relay; the corporate firewall sees ordinary outbound 443 traffic; no public IP, no inbound rule, no port forwarding ever touches the database network segment.

\begin{pitfall}
Installing the gateway on a random analyst workstation is the classic anti-pattern: the machine sleeps, gets patched at whim, and leaves the estate when the analyst does. Gateway hosts are servers---dedicated, patched, monitored, and clustered.
\end{pitfall}

\section{Tuesday: Connecting Hybrid Sources}
\index{SQL Server!gateway}\index{Oracle!gateway}\index{SAP!gateway}
\subsection{SQL Server}
The reference case. Create a connection in Fabric bound to the gateway cluster, supplying server, database, and a least-privilege SQL identity. Pipelines, dataflows, and semantic models then reference the connection by name; the gateway handles the wire. For extraction to OneLake, pair the connection with Chapter~9's metadata-driven Copy pattern.

\subsection{Oracle and SAP}
Oracle connects through the same gateway with the Oracle client or managed provider installed on the gateway host---which makes the host's software inventory part of the integration contract. SAP arrives through the SAP connectors (for example SAP HANA, SAP BW, or SAP Table/BW Open Hub patterns), typically with additional licensing and driver considerations on the gateway host. In both cases, the driver's version and bitness on the host are first-class deployment dependencies.

\subsection{Connection Hygiene}
\begin{itemize}
    \item One logical connection per source per environment (dev/test/prod), named consistently.
    \item Least-privilege identities: read-only accounts for extraction, scoped to required schemas.
    \item Credential rotation procedures documented per connection; a gateway connection with a dead password pages someone at 3 a.m.
\end{itemize}

\begin{tipbox}
Keep a connection inventory table---source, environment, gateway cluster, identity, owner, rotation date---in the same config database as Chapter~9's ingestion config. Connections are configuration, and configuration belongs in tables, not in portal archaeology.
\end{tipbox}

\section{Wednesday: REST APIs and Webhooks}
\index{REST API ingestion}\index{webhooks}\index{pagination}
Not every source is a database. REST APIs and webhooks move an increasing share of enterprise data, and both have discipline requirements that naive pipelines ignore.

\subsection{REST APIs}
Ingest APIs with the Copy activity's REST connector or a Web activity followed by processing. The four recurring concerns: \emph{authentication} (OAuth2 client credentials or managed identity---never API keys pasted into pipeline JSON); \emph{pagination} (follow \texttt{next} links or page parameters until exhaustion; a pipeline that reads page one is a pipeline that loses data); \emph{rate limits} (honor \texttt{429} responses with backoff, or the API will do the throttling for you); and \emph{schema drift} (JSON responses evolve silently---apply Week~9's contract gate to the parsed result).

\subsection{Webhooks}
Webhooks invert the flow: the source pushes events to an endpoint you expose. In Fabric, eventstreams (Chapter~13) and custom endpoints receive such pushes. The credential pattern matters: validate shared secrets or tokens on receipt, and design the receiver to be idempotent because push sources retry.

\begin{pitfall}
Webhook receivers without authentication are open write endpoints into your lake. Validate a shared secret or signature on every request, log rejected calls, and rotate the secret on a schedule.
\end{pitfall}

\section{Thursday: Hands-on Project}
\index{hands-on lab}\index{on-premises data gateway!installation}
The Week~11 lab stands up the full hybrid path: a mock on-premises SQL Server, the on-premises data gateway, and a daily Fabric pipeline extracting through it to OneLake.

\subsection{Goal and Architecture}
Use a local SQL Server (Developer edition, a Docker container, or an instructor VM) seeded with a \texttt{sales.orders} table. Install the on-premises data gateway on a suitable host (your machine is acceptable \emph{for the lab}; note explicitly why it would not be in production). Create the Fabric connection, then build pipeline \texttt{PL\_daily\_orders\_extract} running the extraction on a schedule.

\subsection{Step-by-Step Actions}
\begin{enumerate}
    \item Seed the mock SQL Server with \texttt{sales.orders} (a few thousand rows) and create a read-only login \texttt{fabric\_reader}.
    \item Download and install the on-premises data gateway; register it to your tenant with a recovery key stored in the team vault.
    \item In Fabric, create the connection: gateway, server, database, and the \texttt{fabric\_reader} identity.
    \item Build the pipeline: Copy activity from the gateway connection to \texttt{bronze.orders\_onprem} in the lakehouse, partitioned by load date.
    \item Add the watermark pattern: \texttt{WHERE order\_ts > @watermark} with the watermark advanced only on success.
    \item Schedule the pipeline daily; run it manually once and validate row counts.
    \item Simulate failure (stop the gateway service), observe the pipeline failure, restart, and verify the rerun is idempotent.
\end{enumerate}

\begin{lstlisting}[language=SQL,caption={Mock source setup for the Week 11 lab.},label={lst:chapter11-lab-setup}]
CREATE LOGIN fabric_reader WITH PASSWORD = '<from-vault>';
CREATE USER fabric_reader FOR LOGIN fabric_reader;
EXEC sp_addrolemember 'db_datareader', 'fabric_reader';

CREATE TABLE sales.orders (
    order_id     INT IDENTITY PRIMARY KEY,
    customer_id  INT,
    order_ts     DATETIME2,
    net_amount   DECIMAL(18,2)
);
\end{lstlisting}

\subsection{Validation Checklist}
\begin{itemize}
    \item The gateway registers and shows online in the Fabric portal.
    \item The connection uses the read-only identity, never a personal or admin account.
    \item The pipeline extracts incrementally and the watermark advances only on success.
    \item Gateway-down produces a clear pipeline failure; recovery replays without duplicates.
    \item The gateway recovery key is vaulted, and the host appears in the connection inventory.
    \item You can articulate why the same setup on an analyst laptop would be a production incident waiting to happen.
\end{itemize}

\section{Friday: Use Case and Architecture}
\index{hybrid connectivity!bank}\index{security!network perimeter}
A retail bank's policy is absolute: transactional databases never have public endpoints---no public IP, no inbound firewall exceptions, no exceptions. Core banking runs on on-premises SQL Server and Oracle. The analytics program needs daily extracts into Fabric.

\subsection{The Pattern}
A clustered on-premises data gateway---two nodes minimum, in the bank's data-center network segment---initiates outbound-only TLS connections to the Azure relay. Fabric connections bind to the cluster with least-privilege, read-only identities per source. Pipelines extract on schedule to OneLake Bronze; Week~9's contract gates validate; the run ledger records every attempt. No packet ever flows inbound to the bank's network; no database listens publicly; the firewall policy sees only outbound 443.

\begin{figure}[H]
    \centering
    \begin{tikzpicture}[node distance=8mm and 9mm]
        \node[store] (db) at (0,0) {SQL Server /\\Oracle (private)};
        \node[stage, minimum width=2.8cm] (gw) at (4.6,0) {Gateway cluster\\(2+ nodes)};
        \node[stage, minimum width=2.6cm] (relay) at (9.4,0) {Azure relay\\(outbound 443)};
        \node[stage] (pl) at (13.8,0) {Fabric\\pipelines};
        \node[store] (lh) at (17.6,0) {OneLake\\Bronze};

        \draw[arrow] (db) -- (gw);
        \draw[arrow] (gw) -- node[above, font=\scriptsize]{outbound only} (relay);
        \draw[arrow] (relay) -- (pl);
        \draw[arrow] (pl) -- (lh);

        \node[note, text width=14cm, align=center] at (8.8,-1.8) {The bank's perimeter sees only outbound TLS. No inbound rule, no public IP, no exposed listener.};
    \end{tikzpicture}
    \caption{Bank-grade hybrid pattern: outbound-only gateway cluster.}
    \label{fig:chapter11-bank}
\end{figure}

\subsection{Hardening the Pattern}
\begin{itemize}
    \item \textbf{Cluster for availability.} A single-node gateway is a single point of failure for every hybrid pipeline; install two or more nodes for high availability and load balancing \cite{microsoftOnPremGateway}.
    \item \textbf{Patch within one release of current.} Gateway software carries credentials in motion; unpatched gateways are both reliability and security debt.
    \item \textbf{Monitor node health centrally.} Node offline, version drift, and relay connectivity are alertable conditions, not portal trivia.
    \item \textbf{Restrict outbound destinations.} Where policy demands, egress-filter the gateway segment to the documented Azure relay endpoints.
    \item \textbf{Separate duties.} Gateway admins, connection owners, and pipeline developers are distinct roles---Chapter~21 formalizes this.
\end{itemize}

\begin{pitfall}
The gateway breaks mid-pipeline more often than teams expect---host reboots, Windows updates, certificate renewals. The lab's deliberate gateway-down test is not theater: every hybrid pipeline needs a designed answer to ``the gateway vanished halfway,'' and the answer is watermark-safe idempotent rerun.
\end{pitfall}

\section{Operational Guidance for Week 11}
Gateways and connections are the least visible and most critical infrastructure in a hybrid Fabric estate. Inventory them, cluster them, patch them, and monitor them.

\subsection{Performance and Reliability}
Place gateway hosts near the sources (latency dominates extraction throughput). Size hosts for concurrent query load. Prefer incremental extraction to keep gateway sessions short. Test failover by stopping a node during a run; the cluster should absorb it.

\subsection{Security and Governance Checklist}
\begin{itemize}
    \item Gateway clusters of two or more nodes for every production estate; single nodes are lab-only.
    \item Least-privilege, read-only source identities; rotation dates in the connection inventory.
    \item Recovery keys vaulted at registration time, not after the first rebuild.
    \item Webhook endpoints authenticated and idempotent.
    \item Gateway host patching cadence documented and monitored.
    \item No public endpoints on any source database, verified by network scan, not by assumption.
\end{itemize}

\section{Interview Q\&A}
\begin{enumerate}
    \item \textbf{Q: When do you need a VNet data gateway instead of an on-premises gateway?}\\
    A: When sources live inside an Azure VNet (Azure SQL MI, VMs, private-endpoint services) and you want Microsoft-managed gateway compute without maintaining a Windows host.
    \item \textbf{Q: Why cluster gateway nodes for production workloads?}\\
    A: A single node is a single point of failure for every hybrid pipeline; clusters provide high availability during patches and failures plus load distribution across concurrent extractions.
    \item \textbf{Q: Which credential pattern should webhook/API ingestion use?}\\
    A: OAuth2 client credentials or managed identity, stored in managed connections or Key Vault, with webhook receivers validating a shared secret or signature on every call---never static keys embedded in pipeline JSON.
    \item \textbf{Q: How does a gateway keep data inside the bank's network perimeter?}\\
    A: The gateway initiates every connection outbound over TLS to the Azure relay; the bank's firewall sees only outbound 443 traffic, and no inbound rule or public endpoint ever exists on the database segment.
    \item \textbf{Q: What breaks if the gateway machine goes offline mid-pipeline?}\\
    A: The in-flight activity fails; because extraction is watermark-based and the watermark advances only on success, the rerun replays the same slice idempotently---which is why the failure drill is part of the lab.
    \item \textbf{Q: Why does driver inventory on the gateway host matter for Oracle and SAP?}\\
    A: The host executes the provider code; driver version and bitness are deployment dependencies, and an unmanaged host upgrade can break every Oracle or SAP connection simultaneously.
\end{enumerate}

\section{Further Reading}
Review the on-premises data gateway documentation \cite{microsoftOnPremGateway} and the VNet data gateway documentation \cite{microsoftVNetGateway}. Pipeline mechanics come from Chapter~9 \cite{microsoftFabricPipelines}; eventstream-based webhook reception previews Chapter~13 \cite{microsoftFabricEventstreams}. Security principles here are deepened in Chapter~21 \cite{microsoftFabricSecurity}.

\section*{Key Takeaways}
\begin{itemize}
    \item On-premises gateways are customer-managed agents for corporate networks; VNet gateways are Microsoft-managed attachments for Azure VNets.
    \item The outbound-only relay pattern is why gateway-based hybrid connectivity passes bank-grade security review.
    \item Connections are configuration: inventory them, scope them least-privilege, and rotate their credentials on schedule.
    \item API ingestion demands pagination, rate-limit handling, credential discipline, and schema-drift gates; webhooks demand authentication and idempotent receivers.
    \item Single-node gateways are lab conveniences and production liabilities; cluster everything that matters.
    \item Gateway failure mid-pipeline is a designed-for scenario, not an edge case---watermark-safe reruns absorb it.
    \item The gateway host's patch level and driver inventory are part of the integration contract.
\end{itemize}

\section{Exercises}
\begin{enumerate}
    \item \textbf{(Conceptual)} Explain the outbound-only relay pattern to a network engineer skeptical of ``cloud access to our database.''
    \item \textbf{(Conceptual)} Contrast the on-premises and VNet gateways across management burden, reachable sources, and failure domains.
    \item \textbf{(Hands-on)} Add a second gateway node, form a cluster, and stop one node during an extraction; document the observed behavior.
    \item \textbf{(Hands-on)} Extend the lab pipeline to call a paginated REST API, following \texttt{next} links to exhaustion.
    \item \textbf{(Design)} Design the connection inventory schema: columns, ownership, rotation tracking, and the query that lists connections overdue for rotation.
    \item \textbf{(Design)} The bank acquires a subsidiary whose sources sit in AWS. Propose the connectivity pattern and name every trust boundary crossed.
    \item \textbf{(Interview)} In five sentences, defend gateway clustering as an availability requirement rather than an optimization.
    \item \textbf{(Operational)} Write the runbook for ``gateway node reports unhealthy''---detection, failover verification, patching, and rejoin.
\end{enumerate}

\section{Capstone Project: Clustered Hybrid Extraction with Failover Evidence}\label{sec:ch11-capstone}
\index{capstone project}\index{on-premises data gateway!capstone}\index{hybrid connectivity!capstone}

\begin{capstonebox}
\textbf{Mission.} Build the bank-grade hybrid pattern end to end: a two-node gateway cluster fronting a mock on-premises SQL Server, a scheduled incremental extraction to OneLake, a connection inventory, and documented failover evidence proving the estate survives a node loss mid-run.

\textbf{Estimated time.} 5--7 hours (a second node may be a VM or a second machine; document your topology either way).

\textbf{What you'll practice.}
\begin{itemize}
    \item Registering and clustering gateway nodes with vaulted recovery keys.
    \item Building watermark-safe incremental extraction through the cluster.
    \item Executing and documenting a mid-run failover drill.
    \item Operating connections as inventoried, least-privilege configuration.
\end{itemize}
\end{capstonebox}

\subsection*{Scenario}
Contoso Bank's analytics platform must extract nightly from an on-premises core system under a no-public-endpoint policy. The platform team must demonstrate---to the security review board, with evidence---that extraction is outbound-only, survives a gateway node failure, and never double-counts after a mid-run outage.

\subsection*{Requirements}
\begin{itemize}
    \item Two-node gateway cluster registered to the tenant; recovery keys vaulted.
    \item Least-privilege read-only source identity; connection recorded in the inventory table.
    \item Pipeline \texttt{PL\_nightly\_core\_extract}: watermark-incremental, success-only advance, ledger row per run.
    \item Failover drill: stop one node mid-run, capture the outcome, and prove the retry is duplicate-free.
    \item A network-evidence appendix: firewall or connection logs showing outbound-only flows.
    \item Security review narrative suitable for the bank's review board.
\end{itemize}

\subsection*{Implementation Steps}
\begin{enumerate}
    \item Install and register both gateway nodes; form the cluster; vault the recovery keys.
    \item Create the least-privilege identity and the Fabric connection bound to the cluster.
    \item Build and schedule the incremental extraction pipeline with ledger logging.
    \item Run the failover drill; capture pipeline and gateway evidence.
    \item Assemble the network evidence and the security narrative.
\end{enumerate}

\subsection*{Deliverables}
\begin{itemize}
    \item Cluster topology diagram and node inventory.
    \item Pipeline export and ledger contents across the drill.
    \item Failover evidence: timestamps, affected run, clean retry.
    \item Network evidence and the review-board narrative.
\end{itemize}

\subsection*{Acceptance Criteria}
\begin{itemize}
    \item[\(\square\)] Both nodes serve the cluster; stopping either one does not lose a night's data.
    \item[\(\square\)] The mid-run failure produces a visible failure and a duplicate-free retry.
    \item[\(\square\)] The source identity is read-only and schema-scoped.
    \item[\(\square\)] Network evidence shows outbound-only flows; no inbound rule exists.
    \item[\(\square\)] The connection inventory row is complete, including rotation date.
    \item[\(\square\)] No secrets appear in the repository, pipeline JSON, or screenshots.
\end{itemize}

\subsection*{Stretch Goals}
\begin{itemize}
    \item Add node-health alerting so the drill's failure would have paged the on-call.
    \item Add a second source (Oracle XE or a container) to the same cluster and document the driver dependency.
    \item Automate the connection-inventory freshness check as a scheduled notebook.
    \item Extend the security narrative with an egress-filtering design for the gateway segment.
\end{itemize}
